% Mini-Project 2 - Mid Submission: peer feedback log (half a page or less)
% Build: tectonic feedback.tex   (or: pdflatex feedback.tex)

\documentclass[11pt]{article}
\usepackage[margin=1in]{geometry}
\usepackage[T1]{fontenc}
\usepackage{tgtermes}
\usepackage[hidelinks]{hyperref}

\setlength{\parindent}{0pt}
\setlength{\parskip}{0.6em}
\pagestyle{empty}

\begin{document}

%--------------------------------------------------------------------
% Cover page (same layout as the proposal)
\begin{titlepage}
\centering
\vspace*{1.5in}
{\Huge\bfseries Peer Feedback Log\par}
\vspace{0.25in}
{\Large Mini-Project 2: Mid Submission (Working Prototype)\par}
\vspace{0.1in}
{\large What Drives Medical Insurance Charges?\par}
\vspace{1in}

{\large
\renewcommand{\arraystretch}{1.6}
\begin{tabular}{@{} r l @{}}
\textbf{Name:} & Ashik M John \\
\textbf{Course:} & Information Visualization (CS-GY 6313), Section 1-A \\
\textbf{Instructor:} & Saremi Raz \\
\textbf{Date of Submission:} & 12 November 2026 \\
\end{tabular}
\par}

\vfill

\begin{minipage}{0.85\textwidth}
\raggedright
\textbf{Use of AI Tools.} An AI coding assistant was used for structure and formatting only. Every decision was made by the author.
\end{minipage}
\end{titlepage}

{\large\textbf{Peer Feedback Log: Mid Submission (Working Prototype)}}\\[0.2em]
\textbf{Project:} What drives medical insurance charges? \quad
\textbf{Peer:} Perumal Marimuthu

\textbf{Feedback I received.} Perumal reviewed my dashboard and said it has a strong prototype foundation. He said the D3 scatterplot and bar-chart views are clearly connected, and that the brushing, bar selection, smoker filter, reset option and tooltips make the interaction meaningful rather than just decorative. He also said my \texttt{CHANNELS.md} explains the encoding choices well, especially position for charges and age, hue for smoker status, and ordered placement for BMI categories. The biggest improvement he recommended before the final is a short, visible data-source and preparation note inside the dashboard itself, not only in the README. He suggested it should say the dataset has 1,338 records, name the Kaggle source, and briefly explain the derived BMI and age categories. This would make the data preparation clear during a live presentation, and viewers would not need to open the repository documentation to understand the context. For the final, I plan to make that note easier to see in the dashboard and expand it into a short ``About the data'' section.

\textbf{Feedback I gave to Perumal (Heat Signal).} I told Perumal that Heat Signal is a well-prepared dashboard: the full 264,750-record NYC 311 extract loads and \texttt{prepare\_data.py} checks the record total, the four views each answer a stated question, the cross filtering works (in his Brooklyn screenshot the filter chip, KPI total and monthly chart update together), and the note that 311 reports are signals. My one suggestion was to replace the teal-to-yellow-to-orange heatmap ramp with a single-hue, light-to-dark ramp, because count is an ordered quantity and a mixed-hue ramp is hard to rank and hard to read for colourblind viewers or in greyscale.

\textbf{Planned changes for the final.} (1) In response to Perumal's feedback, I will add a visible ``About the data'' section inside the dashboard, larger and higher on the page than the current small footer, giving the Kaggle source, the 1,338 records, and how the BMI category and age group are derived. (2) I will run a greyscale and colourblind check on the final screenshots to confirm the orange and blue smoker colours stay distinguishable. (3) I will run usability tests with two peers on three tasks each and write up the findings. (4) I will polish the layout and update the proposal and channel table to say D3 and to include the smoker toggle.

\end{document}
